% =============================================================================
% ISI Country Brief — Romania
% External Supplier Concentration Profile
%
% ISI Paper Series — Country Brief
% International Sovereignty Institute — Research Unit
% March 2026
% =============================================================================
\documentclass[12pt,a4paper]{article}

% --- ISI Paper Series shared template ----------------------------------------
\usepackage{isi-series}

% --- PDF metadata ------------------------------------------------------------
\hypersetup{
  pdftitle={External Supplier Concentration Profile: Romania — ISI EU-27 2024},
  pdfauthor={International Sovereignty Institute -- Research Unit},
  pdfcreator={International Sovereignty Institute},
  pdfsubject={ISI Paper Series -- Country Brief RO, EU-27, Vintage 2024},
  pdfkeywords={sovereignty index, EU-27, supplier concentration,
    romania, country brief},
}

% ==============================================================================
\begin{document}
\hypersetup{pageanchor=false}

% ---------- Title Page with Metadata ------------------------------------------
\begin{titlepage}
\centering
\vspace*{2cm}
{\Large\textsc{ISI Paper Series --- Country Brief}}\\[1.2cm]
{\LARGE\bfseries External Supplier Concentration Profile:\\[0.3em]
Romania}\\[1.5cm]
{\large International Sovereignty Index (ISI)\\[0.3em]
EU-27 --- Vintage 2024 --- Methodology v1.0}\\[1.2cm]
{\large International Sovereignty Institute --- Research Unit}\\[0.8cm]
{\large March 2026}

\vspace{0.8cm}
\noindent\rule{\textwidth}{0.4pt}
\vspace{0.4cm}

% --- Research Metadata --------------------------------------------------------
\begin{flushleft}
\noindent\textbf{Data basis:} ISI Paper~2 --- EU-27 Empirical Results 2024 (Methodology~v1.0, Vintage~2024).\\[4pt]
\noindent\textbf{Methodology:} ISI Paper~1 --- Methodological Foundations, Version~1.0.

\medskip

\noindent\textbf{JEL Codes:} F02, F15, F52, O30, Q43\\[4pt]
\noindent\textbf{Keywords:} sovereignty index; EU-27; supplier concentration; composite indicator; variance decomposition; defence dependence

\medskip

\noindent\textbf{Related publications in the ISI Paper Series:}
\begin{itemize}[nosep,leftmargin=1.5em]
\item ISI Paper~1 --- Technical Specification (v1.0) --- doi:\\
  \url{https://doi.org/10.5281/zenodo.18764227}
\item ISI Paper~2 --- EU-27 Empirical Results 2024 (v1.0) --- doi:\\
  \url{https://doi.org/10.5281/zenodo.18764170}
\end{itemize}
\end{flushleft}
\vfill
\end{titlepage}

\hypersetup{pageanchor=true}
\pagenumbering{arabic}
\pagestyle{fancy}
\fancyhead[L]{\small\sffamily\textit{ISI Country Brief}}
\fancyhead[R]{\small\sffamily\textit{Romania}}
\renewcommand{\headrulewidth}{0.4pt}

% ---------- Body --------------------------------------------------------------

\section{Executive Summary}
\label{sec:summary}

Romania ranks~18th in the EU-27 with a composite \ISI score of~0.319, classified as \emph{moderately concentrated}.  The dominant axes are defence~(0.564) and energy~(0.406).  The structural profile is characterised as defence-led.  Under Dirichlet weight perturbation (10\,000~Monte Carlo draws), Romania's mean rank is~17.43 with a standard deviation of~1.65; the 5th--95th percentile band spans ranks~15 to~19.


\section{Country Position in the EU-27}
\label{sec:position}

Romania occupies rank~18 of~27 EU member states (composite score~0.319).  The EU-27 mean composite score is~0.344 with a standard deviation of~0.070.  Romania's score lies below the EU-27 mean, indicating below-average external supplier concentration.  In the ranking, Romania is situated between Portugal (rank~17; 0.322) and Czechia (rank~19; 0.318).


\section{Axis Profile}
\label{sec:axis-profile}

\Cref{tab:axis-profile} presents Romania's scores on all six \ISI axes.

\begin{table}[htbp]
\centering\small
\caption{Romania's \ISI axis profile.}
\label{tab:axis-profile}
\begin{tabular}{@{}lr@{}}
\toprule
\textbf{Axis} & \textbf{Score} \\
\midrule
    Finance           & 0.200 \\
    Energy            & 0.406 \\
    Technology        & 0.297 \\
    Defence           & 0.564 \\
    Critical Inputs   & 0.158 \\
    Logistics         & 0.288 \\
\bottomrule
\end{tabular}
\end{table}

The highest axis score is defence~(0.564), which lies below the EU-27 axis mean of~0.573.  The second-highest axis is energy~(0.406), above the EU-27 axis mean of~0.365.  The lowest axis score is critical inputs~(0.158), below the EU-27 mean of~0.252.


\section{Structural Drivers}
\label{sec:drivers}

\begin{sloppypar}
Romania's composite score is led by defence~(0.564), with energy~(0.406) as the second contributor.  The gap between the two highest axes is~0.158.  While defence exerts the strongest influence on the composite, the contribution of energy is non-negligible.
\end{sloppypar}


\section{Comparative Context}
\label{sec:comparative}

Across the EU-27, the covariance-based variance decomposition shows that defence (35.1\,\%) and logistics (34.3\,\%) together account for 69.4\,\% of total cross-sectional composite variance.  Romania's defence score~(0.564) lies below the EU-27 defence mean~(0.573), and its logistics score~(0.288) lies below the EU-27 logistics mean~(0.518).  Under Dirichlet weight perturbation, Romania's rank standard deviation is~1.65, indicating moderate rank stability.



Romania's profile shares the defence-led structure common across much of the EU-27, but with a distinctive secondary emphasis on technology~(0.297)---the third-highest axis, exceeding the EU-27 technology mean of~0.211.  Energy~(0.406) also lies above its EU-27 mean.  By contrast, logistics~(0.288) and critical inputs~(0.158) are well below their respective averages, producing a profile in which the two highest-variance axes pull in opposite directions.

This cross-cutting axis structure is reflected in the Dirichlet results.  The rank standard deviation of~1.65 and the percentile band~(15--19) are moderate, but the band is positioned entirely in the lower half of the distribution.  Unlike countries with high defence and high logistics, Romania's below-mean logistics score prevents the composite from benefiting fully from the defence elevation.

The composite gap to Portugal~(rank~17; 0.003) and Czechia~(rank~19; 0.001) is small in both directions.  Romania occupies a position in the distribution where three to four countries are separated by less than 0.01 composite-score points---a density that renders ordinal comparisons within this band analytically uninformative without reference to the underlying axis compositions.


\section{Interpretation Boundary}
\label{sec:interpretation}

The \ISI measures concentration of external suppliers, not sovereignty in a normative or political sense.  High concentration may reflect geography, economic specialisation, or alliance structures.  A high composite score does not imply policy failure; a low score does not imply policy success.  The index provides an empirical measurement basis --- what policymakers derive from it is a separate question.


% ---------- References --------------------------------------------------------
\clearpage
\vspace*{1.5cm}

\begingroup\emergencystretch=2em\hbadness=10000
\section*{References}

\begin{enumerate}[label={[\arabic*]},leftmargin=2em,itemsep=6pt]
\item International Sovereignty Institute.
  \textit{International Sovereignty Index (ISI) technical specification, version~1.0}.
  Zenodo, 2026.\\ISI Paper Series, Paper~1.
  doi:\,\href{https://doi.org/10.5281/zenodo.18764227}{10.5281/zenodo.18764227}.

\item International Sovereignty Institute.
  \textit{International Sovereignty Index (ISI): EU-27 empirical results 2024 (v1.0)}.
  Zenodo, 2026.\\ISI Paper Series, Paper~2.
  doi:\,\href{https://doi.org/10.5281/zenodo.18764170}{10.5281/zenodo.18764170}.
\end{enumerate}
\endgroup

\end{document}
